% --- listing style for problem p1 (namespaced) ---
\definecolor{p1kw}{RGB}{0,80,160}
\definecolor{p1cm}{RGB}{110,110,110}
\definecolor{p1st}{RGB}{160,60,20}
\lstdefinestyle{p1py}{%
  language=Python,
  basicstyle=\ttfamily\scriptsize,
  keywordstyle=\color{p1kw}\bfseries,
  commentstyle=\color{p1cm}\itshape,
  stringstyle=\color{p1st},
  showstringspaces=false,
  breaklines=true, breakatwhitespace=true, breakindent=2em,
  columns=fullflexible, keepspaces=true,
  xleftmargin=0.6em, aboveskip=0.8em, belowskip=0.4em,
}
\section{Q1: Structural Encoding in Graph Transformer}

\texttt{LapPETransformer} adds Laplacian eigenvector positions to every node once, before the
first block; \texttt{SPDAttention} adds a learnable per-head scalar, indexed by the pair's hop
distance, to every attention score. Two choices in the absolute encoding matter.

\textbf{Which eigenvectors, and at what scale.} The smallest eigenvalue of
$L = I - D^{-1/2} A D^{-1/2}$ is $0$ with the trivial eigenvector $D^{1/2}\mathbf{1}$, so
\texttt{precompute} drops it and returns columns
$1,\dots,\texttt{LAP\_K}$; a graph with fewer non-trivial eigenvectors is zero-padded (ZINC's
smallest molecule has $9$ nodes, so exactly $8$ exist). Each column is then scaled by $\sqrt{n}$:
eigenvector entries are $O(n^{-1/2})$, so unscaled a channel has RMS $0.011$ on the $7624$-node
LastFM graph, where \texttt{num\_atom\_types=0} leaves the encoding as the model's only input.

\textbf{Sign ambiguity.} $v$ and $-v$ are the same eigenvector and the solver picks one
arbitrarily, so each graph's channels get an independent random sign flip while training and the
identity at evaluation. That is an augmentation \emph{across} graphs, so it is skipped when the
batch holds only one --- on LastFM it would be noise on the model's only input.

The relative encoding computes all-pairs hop distances by frontier expansion: exact on ZINC,
blocked and stopped after \texttt{max\_dist} hops on LastFM; distances from \texttt{max\_dist}
upwards share one bias anyway, and unreachable pairs keep their own row. The lookup is added to
the scores before mask and softmax,
$\mathrm{score}_{ij} = q_i^\top k_j/\sqrt{d} + b_{\mathrm{SPD}(i,j)}$, which is Graphormer's
form. Distances were checked against a plain BFS, eigenvectors against the dense spectrum of
$L$.

\begin{table}[H]
\centering
\caption{ZINC subset, 30 epochs, dense attention; 4 blocks, dim 64, 4 heads, the framework's fixed seed. Epoch selected on validation MAE.}
\label{p1:tab:q1}
\begin{tabular}{lccc}
\toprule
encoding & val MAE & test MAE & best epoch \\
\midrule
none (plain \texttt{Attention}) & 0.6472 & 0.7013 & 18 \\
demo: degree centrality $+$ common neighbours & 0.4629 & 0.4957 & 27 \\
LapPE only (absolute) & 0.6205 & 0.6575 & 30 \\
SPD only (relative) & 0.5373 & 0.5614 & 29 \\
\textbf{LapPE $+$ SPD} (cell 14) & \textbf{0.5160} & \textbf{0.5496} & 30 \\
\bottomrule
\end{tabular}
\end{table}

SPD carries almost all of the gain over plain attention (0.7013 to 0.5614 test MAE); LapPE alone reaches 0.6575. Putting LapPE on top of SPD moves it by 0.0119, inside the 0.0200 seed spread of the full-attention row of Table~\ref{p1:tab:q2}, so this budget does not establish a LapPE effect beyond SPD. Both choices were ablated on both datasets, and behave differently. The $\sqrt{n}$ scaling is invisible on ZINC (0.5449 against 0.5496) and worth 0.0508 micro-F1 on LastFM. The sign flip is worth 0.0503 test MAE on ZINC and costs 0.1281 micro-F1 on LastFM, which is why it is conditional.

\subsection*{Implementation (notebook cell 13, with its two helpers)}

\begin{lstlisting}[style=p1py]
SPD_MAX_DIST = 8          # distances above this share one bias (SPDAttention.max_dist)
DENSE_LIMIT = 1024        # graphs this small are handled with dense linear algebra


def laplacian_eigvecs(edge_index, num_nodes, k=LAP_K):
    """Rows of the k eigenvectors of L = I - D^-1/2 A D^-1/2 with the smallest
    non-trivial eigenvalues, scaled by sqrt(n) so each channel has unit RMS.
    Padded with zero columns when the graph has fewer than k non-trivial eigenvectors."""
    n = num_nodes
    A = sp.coo_matrix((np.ones(edge_index.size(1)), (edge_index[0].numpy(), edge_index[1].numpy())),
                      shape=(n, n)).tocsr()
    dinv = 1.0 / np.sqrt(np.maximum(np.asarray(A.sum(1)).ravel(), 1.0))
    Ahat = sp.diags(dinv) @ A @ sp.diags(dinv)          # = I - L, same eigenvectors
    m = min(k, n - 1)                                   # non-trivial eigenvectors available
    out = np.zeros((n, k), dtype=np.float32)
    if m > 0:
        if n <= DENSE_LIMIT:
            w, V = np.linalg.eigh(np.eye(n) - Ahat.toarray())           # ascending
        else:                                            # smallest of L = largest of I - L
            v0 = np.random.default_rng(RANDOM_SEED).standard_normal(n)   # fixed ARPACK start
            w, V = sp.linalg.eigsh(Ahat, k=m + 1, which="LA", tol=0, v0=v0)
            w, V = 1.0 - w, V
            order = np.argsort(w)
            w, V = w[order], V[:, order]
        out[:, :m] = V[:, 1:m + 1] * math.sqrt(n)        # skip the trivial eigenvector
    return torch.from_numpy(out)


def hop_distances(edge_index, num_nodes, cap=None):
    """All-pairs hop distance by frontier expansion, as int16 [n, n].

    Unreachable pairs (different connected components) are -1.  With `cap` set the
    expansion stops after `cap` hops and every pair that is further away is reported
    as `cap`; those pairs share the max_dist bias anyway."""
    n = num_nodes
    src, dst = edge_index[0].numpy(), edge_index[1].numpy()
    A = sp.coo_matrix((np.ones(len(src), dtype=np.float32), (src, dst)), shape=(n, n)).tocsr()
    D = np.full((n, n), -1, dtype=np.int16)
    limit = n if cap is None else cap
    if n <= DENSE_LIMIT:
        Ad = np.zeros((n, n), dtype=np.float32)
        Ad[src, dst] = 1.0
        seen = np.eye(n, dtype=bool)
        frontier = seen.copy()
        D[seen] = 0
        for d in range(1, limit + 1):
            frontier = ((frontier.astype(np.float32) @ Ad) > 0) & ~seen
            if not frontier.any():
                break
            D[frontier] = d
            seen |= frontier
    else:
        for s0 in range(0, n, 1024):                     # blocks of sources, columns of F
            s1 = min(s0 + 1024, n)
            F = np.zeros((n, s1 - s0), dtype=np.float32)
            F[np.arange(s0, s1), np.arange(s1 - s0)] = 1.0
            seen = F.astype(bool)
            D[np.arange(s0, s1), np.arange(s0, s1)] = 0
            for d in range(1, limit + 1):
                F = ((A @ F) > 0) & ~seen
                if not F.any():
                    break
                D[s0:s1][F.T] = d
                seen |= F
                F = F.astype(np.float32)
    if cap is not None:                                  # reached later than the cap, not unreachable
        comp = sp.csgraph.connected_components(A, directed=False)[1]
        far = (D < 0) & (comp[:, None] == comp[None, :])
        D[far] = cap
    return torch.from_numpy(D)
\end{lstlisting}

\begin{lstlisting}[style=p1py]
class LapPETransformer(GraphTransformer):
    """Absolute encoding: Laplacian eigenvectors, added to the input once."""

    def __init__(self, *args, **kwargs):
        super().__init__(*args, **kwargs)
        # Linear map from the LAP_K eigenvector entries of a node to a node vector of size dim.
        self.lap_proj = nn.Linear(LAP_K, self.dim)

    @staticmethod
    def precompute(edge_index, num_nodes):
        return {"lap": laplacian_eigvecs(edge_index, num_nodes)}

    def inject_absolute(self, h, batch):
        lap = batch["lap"]
        # Eigenvectors are defined up to sign, so flip each graph's channels at random while
        # training and use the identity at evaluation.  This is an augmentation ACROSS graphs:
        # on a single-graph batch there is nothing to generalise over and the flip is pure
        # input noise, so it is skipped (it costs ~0.12 micro-F1 on LastFM).
        if self.training and batch["num_graphs"] > 1:
            sign = torch.randint(0, 2, (batch["num_graphs"], lap.size(1)),
                                 device=lap.device, dtype=lap.dtype) * 2 - 1
            lap = lap * sign[batch["batch"]]
        return h + self.lap_proj(lap)


class SPDAttention(Attention):
    """Relative encoding: shortest-path distance, one learnable bias per head."""

    def __init__(self, dim, heads, max_dist=SPD_MAX_DIST):
        super().__init__(dim, heads)
        self.max_dist = max_dist
        # one bias per head for each distance 0..max_dist, plus one for unreachable pairs
        self.spd_table = nn.Embedding(max_dist + 2, heads)

    @staticmethod
    def precompute(edge_index, num_nodes):
        cap = None if num_nodes <= DENSE_LIMIT else SPD_MAX_DIST
        return {"spd": hop_distances(edge_index, num_nodes, cap=cap)}

    def bucket(self, d):
        """Hop distance -> row of spd_table: 0..max_dist, and max_dist + 1 if unreachable."""
        d = d.long()
        return torch.where(d < 0, torch.full_like(d, self.max_dist + 1), d.clamp(max=self.max_dist))

    def forward(self, h, enc, mask):
        B, N, D = h.shape
        q, k, v = self.qkv(h).view(B, N, 3, self.heads, D // self.heads).permute(2, 0, 3, 1, 4)
        scores = q @ k.transpose(-1, -2) / math.sqrt(D // self.heads)      # [B, H, N, N]
        bias = self.spd_table(self.bucket(enc["spd"]))                     # [B, N, N, H]
        scores = scores + bias.permute(0, 3, 1, 2)
        scores = scores.masked_fill(~mask[:, None, None, :], float("-inf"))
        weights = torch.softmax(scores, dim=-1).nan_to_num(0.0)
        return self.out((weights @ v).transpose(1, 2).reshape(B, N, D))
\end{lstlisting}

\section{Q2: Efficient Graph Transformer}

\texttt{SparseSPDAttention} keeps Q1's parameters and \texttt{precompute} and changes only which
pairs exist. Five pair sets were compared (Table~\ref{p1:tab:q2}): the graph's edges
alone; plus one global virtual node; plus Exphormer's random expander at two widths; and the
2-hop ball --- every ordered pair within two hops --- plus the virtual node. The
submitted class attends to the last of these, picked on validation. Its cost is
$\sum_v \deg(v)^2$ rather than $n^2$: 189 pairs per ZINC molecule, 804{,}349
on LastFM ($1.4\%$ of full). The expander was the a-priori favourite: it spreads global
mixing over random edges instead of routing it through one hub, still at $O(n)$ pairs. Both
widths finished last on both datasets.

\texttt{forward} gathers $q,k,v$ at the endpoints of each pair, adds the inherited SPD bias plus
one learnable bias per pair type, normalises with \texttt{segment\_softmax} over the pairs that
share a querying node, and scatters the weighted values back with \texttt{index\_\allowbreak{}add\_}. The
handout points at a \texttt{segment\_\allowbreak{}attention} helper that does not exist in PyG 2.8.1; cell 3
imports \texttt{softmax as\allowbreak{} segment\_\allowbreak{}softmax}, and that is what is used. The framework also
reports \texttt{enc["spd"]}$=0$ for pairs touching a virtual node, which would alias them with
distance-$0$ self loops, so those go to the unreachable row of \texttt{spd\_table} instead. On a complete pair set this path reproduces the dense Q1 attention to
$7.5\times10^{-8}$.

\begin{table}[H]
\centering
\caption{Masking strategies on ZINC, same model and budget as Table~\ref{p1:tab:q1}, with LapPE throughout. Pairs are per graph, averaged over the 10{,}000 training molecules ($n^2 = 557$ on average). The val and test columns are the framework's seed, as cell 18 runs it; the last column is the mean and half-range of the test MAE over that seed and two more, the expander rows redrawing their expander each time. The bolded row has not been reproduced exactly, in either column: three later repeats of that identical configuration --- same seed, same weights, same batch order --- gave val MAE 0.4318, 0.4418, 0.4418 and test MAE 0.4715, 0.4766, 0.4766, every one of them worse than the draw printed here, because the sparse forward accumulates with \texttt{index\_add\_}, whose order across CPU threads is not fixed. Read any single ZINC number in this table as $\pm$0.01, including the validation column the choice of pair set rests on --- that choice is made on the per-seed ranking, which is stable, not on these decimals. All of this compares pair sets inside this framework --- 30 epochs, and the bond features \texttt{make\_batch} carries are never read --- rather than reproducing any paper's numbers.}
\label{p1:tab:q2}
\small
\begin{tabular}{lrrccc}
\toprule
attention pairs & pairs/graph & of full & val MAE & test MAE & 3-seed mean \\
\midrule
self loops $+$ graph edges & 73 & 13\% & 0.4918 & 0.5048 & 0.5145 $\pm$ 0.0133 \\
$+$ virtual node & 120 & 22\% & 0.4791 & 0.5047 & 0.5092 $\pm$ 0.0067 \\
\textbf{2-hop ball $+$ virtual} (submitted) & 189 & 34\% & \textbf{0.4293} & \textbf{0.4677} & 0.4909 $\pm$ 0.0211 \\
$+$ 4-regular expander & 200 & 36\% & 0.5179 & 0.5551 & 0.5708 $\pm$ 0.0129 \\
$+$ 8-regular expander & 266 & 48\% & 0.5213 & 0.5554 & 0.5655 $\pm$ 0.0119 \\
\midrule
full attention (Q1) & 557 & 100\% & 0.5160 & 0.5496 & 0.5614 $\pm$ 0.0100 \\
\bottomrule
\end{tabular}
\end{table}

The ranking is not the one the pair counts suggest. Restricting attention to the graph beats full attention outright --- 0.5145 against 0.5614 test MAE on the three-seed mean --- so on 23-atom molecules sparsity is a useful prior, not a sacrifice. Two hops wins again: lowest validation MAE of the five at every seed, and 0.0184 off the runner-up's three-seed test mean, while a comparable random expander loses 0.0799 on it: the pairs have to mean something, and an expander edge means nothing to a molecule. LastFM neither contradicts nor confirms that --- its top three finish within 0.008 validation micro-F1, the submitted set second of five --- so ZINC decides. Table~\ref{p1:tab:q2lastfm} is what LastFM does settle, the cost: 0.33\,s a step against 5.23\,s for full attention over all 58{,}125{,}376 pairs, a 16$\times$ saving, at test micro-F1 0.7827.

\begin{table}[H]
\centering
\caption{LastFM (7624 nodes, 27{,}806 edges, 18 countries), 2 blocks, dim 32, 2 heads, CPU, 8 threads. One training step is the best of 7 forward-backward passes; micro-F1 is \texttt{run\_node\_task} over 200 epochs, selected on validation. The framework's own single-step \texttt{benchmark\_step} prints 5.24\,s and 0.33\,s for the two models of cell 19, the same numbers. Full attention is not trained for 200 epochs: at 5.23\,s a step it would cost more than every other experiment here together. The framework's peak-memory column prints \texttt{n/a on CPU}, so no memory figure is quoted.}
\label{p1:tab:q2lastfm}
\small
\begin{tabular}{lrrrc}
\toprule
attention pairs & pairs attended & of full & one step (s) & test micro-F1 \\
\midrule
full attention (Q1) & 58{,}125{,}376 & 100.000\% & 5.23 & --- \\
self loops $+$ graph edges & 63{,}236 & 0.109\% & 0.03 & 0.7893 \\
$+$ virtual node & 78{,}485 & 0.135\% & 0.04 & 0.7826 \\
\textbf{2-hop ball $+$ virtual} (submitted) & 804{,}349 & 1.384\% & 0.33 & \textbf{0.7827} \\
$+$ 4-regular expander & 108{,}965 & 0.187\% & 0.05 & 0.7678 \\
$+$ 8-regular expander & 139{,}417 & 0.240\% & 0.06 & 0.7607 \\
\bottomrule
\end{tabular}
\end{table}

\subsection*{Implementation (notebook cell 16)}

\begin{lstlisting}[style=p1py]
EXPANDER_CYCLES = 2       # union of 2 random Hamiltonian cycles = a 4-regular expander


def pair_set(edge_index, num_nodes, n_cycles=EXPANDER_CYCLES, virtual=True, extra=()):
    """The pair set of one graph: self loops, the graph's edges, an expander built as the
    union of n_cycles random Hamiltonian cycles, and one global virtual node (id num_nodes).
    The cycles are drawn from a generator seeded by the graph itself, so they are fixed."""
    n, node = num_nodes, torch.arange(num_nodes)
    gen = torch.Generator().manual_seed((RANDOM_SEED * 1000003 + n * 7919
                                         + int(edge_index.sum())) % (2 ** 31))
    parts = [torch.stack([node, node]), edge_index, *extra]        # self loops + graph edges
    for _ in range(n_cycles if n >= 3 else 0):                     # expander edges
        p = torch.randperm(n, generator=gen)
        q = torch.roll(p, -1)
        parts += [torch.stack([p, q]), torch.stack([q, p])]
    if virtual:                                                    # global node, both directions
        virt = torch.full_like(node, n)
        parts += [torch.stack([node, virt]), torch.stack([virt, node]), torch.tensor([[n], [n]])]
    return coalesce(torch.cat(parts, 1), num_nodes=n + 1)


class SparseSPDAttention(SPDAttention):
    """Your Q1 attention, restricted to a chosen set of node pairs.

    Of the five pair sets measured (see the report) the 2-hop ball plus one global virtual node
    has the lowest validation MAE on ZINC at every seed; on LastFM it is second of five, 0.005
    behind the best.  That is what this class attends to.  pair_set above still builds
    Exphormer's random expander; the expander variants are one argument away and are compared
    in the report."""

    def __init__(self, dim, heads, max_dist=SPD_MAX_DIST):
        super().__init__(dim, heads, max_dist)
        # Exphormer gives each kind of attention edge its own bias: local, expander, global
        self.type_table = nn.Embedding(len(PAIR_TYPES), heads)

    @staticmethod
    def attention_pairs(edge_index, num_nodes):
        """Every ordered pair at most two hops apart, plus one global virtual node."""
        n = num_nodes
        A = sp.coo_matrix((np.ones(edge_index.size(1)), (edge_index[0].numpy(), edge_index[1].numpy())),
                          shape=(n, n)).tocsr()
        ball = ((A + A @ A) > 0).tocoo()                   # pairs at distance 0, 1 or 2
        extra = torch.from_numpy(np.stack([ball.row, ball.col])).long()
        return pair_set(edge_index, num_nodes, n_cycles=0, virtual=True, extra=[extra])

    def forward(self, h, enc, pairs):
        M, D = h.shape
        H, dh = self.heads, D // self.heads
        src, dst = pairs                                           # node dst attends to node src
        q, k, v = self.qkv(h).view(M, 3, H, dh).permute(1, 0, 2, 3)
        score = (q[dst] * k[src]).sum(-1) / math.sqrt(dh)          # [P, H]
        virtual = enc["pair_type"] == PAIR_TYPES.index("virtual")  # their spd arrives as 0,
        d = torch.where(virtual, torch.full_like(enc["spd"], -1), enc["spd"])  # but no distance
        score = score + self.spd_table(self.bucket(d)) + self.type_table(enc["pair_type"])
        alpha = segment_softmax(score, dst, num_nodes=M)           # softmax over each node's pairs
        out = torch.zeros(M, H, dh, device=h.device, dtype=h.dtype)
        out.index_add_(0, dst, alpha.unsqueeze(-1) * v[src])       # [M, H, dh]
        return self.out(out.reshape(M, D))
\end{lstlisting}
